\chapter{Kurir Kimia dan Transduksi Isyarat}\label{ch:b2:cell-signalling}

Sebutir molekul adrenalin, yang tiba di permukaan sebuah sel hati,
tidak pernah masuk ke dalamnya. Ia mengikat sebuah protein di luarnya,
mengubah bentuk proteinnya, dan dalam sedetik selnya telah membuat
sepuluh ribu molekul \omterm{prop:b2:cell-signalling:gpcr}{kurir kedua}, mengaktifkan seratus ribu molekul
enzim, lalu mulai melepaskan sejuta molekul glukosa --- sebuah rantai
penguatan yang panjangnya sejuta kali lipat, yang dimulai oleh sebuah
isyarat yang tetap tinggal di luar pintunya. Bab ini membahas kurir
yang dikirim sel satu sama lain dan bagaimana sebuah sel mengubah
sebuah pesan di permukaannya menjadi sebuah perubahan pada kimianya,
gennya, atau bentuknya: reseptornya, \omterm{prop:b2:cell-signalling:gpcr}{kurir keduanya}, riam kinasenya,
dan kedua rancangannya --- pensinyalan saraf yang cepat dan terkabel
serta pensinyalan \omterm{def:b2:cell-signalling:kinds}{hormon} yang lambat dan tersiar --- yang
pertentangannya berjalan sepanjang sisa buku ini.

\section{Kurir}

\begin{definition}[Macam pensinyalan kimia]\label{def:b2:cell-signalling:kinds}
\index{hormon}\index{pensinyalan endokrin}\index{pensinyalan parakrin}\index{neurotransmiter}
Sebuah sel memberi isyarat kepada sel lain dengan sebuah molekul yang
disekresikan, yaitu sebuah \emph{kurir} atau \emph{ligan}, yang
mengikat sebuah \emph{reseptor} pada atau di dalam sasarannya. Menurut
jangkauannya: pensinyalan \emph{endokrin}, yang di dalamnya sebuah
\emph{hormon} dilepaskan ke dalam \omterm{def:b2:blood-circulation:blood}{darahnya} oleh sebuah kelenjar lalu
mencapai tiap sel tubuhnya, dan hanya bekerja pada sel yang memiliki
reseptornya, selama beberapa menit sampai beberapa jam (insulin,
kortisol, adrenalin, tiroksin); pensinyalan \emph{parakrin}, yang di
dalamnya kurirnya berdifusi ke sel tetangganya dalam jarak satu
milimeter lalu dimusnahkan di jalan (faktor pertumbuhan, histamin,
prostaglandin, morfogen pada \cref{ch:b2:vertebrate-development});
\emph{autokrin}, yaitu selnya memberi isyarat kepada dirinya sendiri;
dan \emph{sinaptik}, yang di dalamnya sebuah neuron melepaskan sebuah
\emph{neurotransmiter} melintasi celah \qty{20}{nm} ke atas satu sel
sasaran, dalam satu milidetik (\cref{ch:b2:neurons-synapses}). Menurut
kimianya: kurir yang \emph{larut air} --- turunan asam amino seperti
adrenalin, peptida dan protein seperti insulin --- tidak dapat
menyeberangi membrannya dan bekerja pada reseptor permukaannya; kurir
yang \emph{larut lipid} --- steroid (kortisol, estradiol, testosteron),
tiroksin, asam retinoat, nitrogen monoksida --- menyeberanginya lalu
bekerja pada reseptor di dalamnya. Sistem endokrin dan sistem sarafnya
adalah dua cara sebuah tubuh berisi $10^{13}$ sel menyelaraskan
dirinya: yang satu menyiarkan perlahan kepada semua orang, yang lain
mengabel dengan cepat kepada seseorang.
\end{definition}

\begin{omfigure}
\begin{tikzpicture}[every node/.style={font=\scriptsize}, >=stealth]
  % endocrine
  \begin{scope}
    \draw[thick, fill=omProp!25] (0,1.6) circle (0.35); \node[above] at (0,2.0) {kelenjar};
    \draw[very thick, omThm!60] (-0.5,0.6) -- (4.5,0.6); \draw[very thick, omThm!60] (-0.5,0.2) -- (4.5,0.2); \node[omThm!70!black] at (2,0.4) {darah};
    \foreach \x in {0.6,1.5,2.4,3.3} \fill[omProp] (\x,0.4) circle (0.05);
    \draw[->, omProp] (0,1.25) -- (0,0.65);
    \draw[thick, fill=omDef!20] (3.8,-0.8) circle (0.35); \node[below, align=center] at (3.8,-1.2) {sel sasaran\\dengan reseptor};
    \draw[->, omProp] (3.8,0.15) -- (3.8,-0.4);
    \draw[thick, fill=gray!20] (1.0,-0.8) circle (0.35); \node[below, align=center] at (1.0,-1.2) {tanpa reseptor:\\mengabaikannya};
    \node[align=center] at (2.2,-2.4) {endokrin: hormon di dalam darah,\\seluruh tubuh, menit sampai jam};
  \end{scope}
  % paracrine
  \begin{scope}[xshift=6.6cm]
    \draw[thick, fill=omProp!25] (0.8,0.6) circle (0.35);
    \foreach \a in {0,45,...,315} \fill[omProp] (0.8,0.6) ++(\a:0.6) circle (0.04);
    \draw[thick, fill=omDef!20] (1.9,0.9) circle (0.3); \draw[thick, fill=omDef!20] (1.7,-0.2) circle (0.3); \draw[thick, fill=omDef!20] (-0.2,0.0) circle (0.3);
    \node[align=center] at (0.8,-2.4) {parakrin: berdifusi ke\\tetangganya, dimusnahkan\\dalam satu milimeter};
  \end{scope}
  % synaptic
  \begin{scope}[xshift=11.6cm]
    \draw[very thick, omDef] (0,1.8) -- (0,0.5); \draw[thick, fill=omDef!25] (0,0.3) circle (0.25);
    \draw[thick, fill=omExo!30] (0,-0.5) circle (0.3);
    \foreach \x in {-0.08,0,0.08} \fill[omProp] (\x,-0.05) circle (0.03);
    \node[align=center] at (0.2,-2.4) {sinaptik: satu sasaran,\\\qty{20}{nm},\\satu milidetik};
  \end{scope}
\end{tikzpicture}

{\small Tiga jangkauan pensinyalan kimia: sebuah \omterm{def:b2:cell-signalling:kinds}{hormon} yang disiarkan
\omterm{def:b2:blood-circulation:blood}{darahnya} ke tiap sel yang membawa reseptornya; sebuah faktor parakrin
yang mencapai tetangganya; sebuah \omterm{def:b2:cell-signalling:kinds}{neurotransmiter} yang diantarkan
melintasi sebuah sinapsis ke satu sel.}
\end{omfigure}

\begin{theorem}[Keterisian reseptor]\label{thm:b2:cell-signalling:occupancy}
\index{tetapan disosiasi}\index{keterisian reseptor}
Sebuah reseptor $R$ mengikat ligannya $L$ secara dapat balik, $R + L
\rightleftharpoons RL$, dengan sebuah \emph{\omterm{thm:b2:cell-signalling:occupancy}{tetapan disosiasi}} $K_d =
[R][L]/[RL]$. Pada kesetimbangan bagian reseptor yang terisi adalah
\[
\theta = \frac{[RL]}{[R]_{\text{tot}}} = \frac{[L]}{K_d + [L]},
\]
yaitu hiperbola yang sama dengan kejenuhan sebuah enzim. Nilai $K_d$
adalah kepekatan yang separuh reseptornya terisi padanya dan mengukur
\emph{afinitasnya}: makin rendah, makin erat. Reseptor \omterm{def:b2:cell-signalling:kinds}{hormon} memiliki
$K_d$ dalam jangkauan nanomolar sampai pikomolar, yang diserasikan
dengan kepekatan \omterm{def:b2:cell-signalling:kinds}{hormon} beredar ($10^{-12}$ sampai $10^{-9}$ mol/L):
pada $[L] = K_d/10$ sepersepuluhnya terisi, pada $10K_d$ sembilan
persepuluhnya. Sebuah sel tidak dapat menanggapi sebuah \omterm{def:b2:cell-signalling:kinds}{hormon} yang
reseptornya tidak ia miliki, dan ia menyetel kepekaannya dengan
mengubah jumlah reseptornya, atau --- seperti pada
\cref{ch:b2:reproductive-hormones} --- dengan menarik reseptornya
ketika \omterm{def:b2:cell-signalling:kinds}{hormonnya} hadir terus-menerus.
\end{theorem}
\begin{proof}
Dengan $[R] = [R]_{\text{tot}} - [RL]$, $K_d = ([R]_{\text{tot}} -
[RL])[L]/[RL]$, sehingga $[RL](K_d + [L]) = [R]_{\text{tot}}[L]$, yaitu
rumusnya. Ia adalah ungkapan Michaelis--Menten pada jilid Tahun 1 tanpa
langkah katalitiknya.
\end{proof}

\begin{omfigure}
\begin{tikzpicture}
\begin{axis}[
  width=0.68\linewidth, height=5.2cm,
  axis x line=bottom, axis y line=left,
  xmode=log, xmin=0.01, xmax=100, ymin=0, ymax=1.1,
  xlabel={kepekatan ligan (dalam satuan $K_d$)}, ylabel={bagian reseptor yang terisi},
  xlabel style={font=\small, at={(axis description cs:0.5,-0.13)}, anchor=north},
  ylabel style={font=\small},
  tick label style={font=\small}, xtick={0.01,0.1,1,10,100}, ytick={0,0.5,1},
  clip=false,
]
\addplot[very thick, omDef, domain=0.01:100, samples=100] {x/(1+x)};
\draw[dashed, gray] (axis cs:1,0) -- (axis cs:1,0.5) -- (axis cs:0.01,0.5);
\node[font=\scriptsize, anchor=north west] at (axis cs:1.1,0.48) {$[L] = K_d$: separuh terisi};
\node[font=\scriptsize, anchor=south east] at (axis cs:0.1,0.1) {10\% pada $K_d/10$};
\node[font=\scriptsize, anchor=north west] at (axis cs:10,0.92) {90\% pada $10\,K_d$};
\end{axis}
\end{tikzpicture}

{\small \omterm{thm:b2:cell-signalling:occupancy}{Keterisian reseptor} terhadap kepekatan ligan pada skala
logaritmik: tanggapan sebuah sel berjalan dari sepersepuluh sampai
sembilan persepuluh sepanjang jangkauan kepekatan seratus kali lipat
yang berpusat pada $K_d$.}
\end{omfigure}

\section{Reseptor di permukaannya}

\begin{proposition}[Reseptor berpasangan protein G dan AMP siklik]\label{prop:b2:cell-signalling:gpcr}
\index{protein G}\index{reseptor berpasangan protein G}\index{AMP siklik}%
\index{kurir kedua}\index{protein kinase A}
Keluarga reseptor yang terbesar --- sekitar \num{800} gen pada manusia,
bagi \omterm{def:b2:cell-signalling:kinds}{hormon}, \omterm{def:b2:cell-signalling:kinds}{neurotransmiter}, bau, dan cahaya --- adalah protein yang
menyeberangi membrannya tujuh kali. Sebuah ligan yang mengikat di
luarnya mengubah bentuknya di dalam, dan di sana ia mengaktifkan sebuah
\emph{\omterm{prop:b2:cell-signalling:gpcr}{protein G}}: sebuah sakelar bersubunit tiga yang menukar GDP yang
terikat padanya dengan GTP, membelah, lalu dalam bentuknya yang terikat
GTP mengatur sebuah enzim efektor selama beberapa detik sebelum
menghidrolisis GTP-nya lalu mematikan dirinya sendiri. Efektor
klasiknya adalah \emph{adenilil siklase}, yang membuat \emph{\omterm{prop:b2:cell-signalling:gpcr}{AMP
siklik}} dari ATP: sebuah \emph{\omterm{prop:b2:cell-signalling:gpcr}{kurir kedua}} yang kecil dan dapat
berdifusi, yang kepekatannya naik dalam hitungan detik dari
\qty{0.1}{\micro mol/L} menjadi beberapa mikromolar dan yang
mengaktifkan \emph{\omterm{prop:b2:cell-signalling:gpcr}{protein kinase A}}, yaitu sebuah enzim yang
memfosforilasi --- dan dengan demikian menyalakan atau memadamkan ---
berlusin-lusin protein sasaran. Isyaratnya diakhiri oleh
fosfodiesterase, yang memusnahkan cAMP-nya, dan oleh fosfatase, yang
menyingkirkan fosfatnya. Adrenalin pada sebuah sel hati bekerja begini:
reseptor $\to$ \omterm{prop:b2:cell-signalling:gpcr}{protein G} $\to$ siklase $\to$ cAMP $\to$ kinase A $\to$
fosforilase kinase $\to$ glikogen fosforilase $\to$ glukosa. Glukagon
memakai rute yang sama; \omterm{prop:b2:cell-signalling:gpcr}{protein G} lainnya menghambat siklasenya, atau
mengaktifkan fosfolipase C, yang melepaskan dua \omterm{prop:b2:cell-signalling:gpcr}{kurir kedua} lainnya,
yaitu \emph{inositol trifosfat} --- yang membuka simpanan kalsiumnya
--- dan diasilgliserol.
\end{proposition}
\begin{proof}[Bukti eksperimental]
Sutherland (1957) menunjukkan bahwa adrenalin yang ditambahkan kepada
sebuah homogenat hati melepaskan glukosa hanya jika membrannya hadir,
bahwa membrannya menghasilkan sebuah faktor yang tahan panas yang
seorang diri mengaktifkan fosforilase di dalam pecahan yang bebas
membran, dan bahwa faktornya adalah \omterm{prop:b2:cell-signalling:gpcr}{AMP siklik} --- \omterm{prop:b2:cell-signalling:gpcr}{kurir kedua} yang
pertama, dan peragaan pertama bahwa sebuah \omterm{def:b2:cell-signalling:kinds}{hormon} bekerja di permukaan
selnya lewat sebuah perantara di dalam selnya. Rodbell dan Gilman
(tahun 1970-an) menemukan bahwa siklasenya memerlukan GTP lalu
mengisolasi \omterm{prop:b2:cell-signalling:gpcr}{protein G} yang membawa isyaratnya antara reseptor dan
enzimnya; toksin kolera, yang mengunci \omterm{prop:b2:cell-signalling:gpcr}{protein G-nya} dalam keadaan
menyala, membuat ususnya menyekresikan berliter-liter cairan tiap hari
--- penyakitnya adalah sebuah cacat pensinyalan.
\end{proof}

\begin{omfigure}
\begin{tikzpicture}[every node/.style={font=\scriptsize}, >=stealth]
  % membrane
  \fill[omExo!30] (0,2.0) rectangle (12,2.5);
  \node[anchor=east] at (-0.1,2.25) {membran};
  % receptor
  \draw[thick, fill=omDef!40] (1.0,1.7) rectangle (1.8,2.8); \node[anchor=east] at (0.95,2.65) {reseptor};
  \fill[omProp] (1.4,3.35) circle (0.12); \node[omProp, above] at (1.4,3.45) {adrenalin};
  \draw[->, omProp] (1.4,3.23) -- (1.4,2.85);
  % G protein
  \draw[thick, fill=omDef!25] (2.6,1.5) ellipse (0.72 and 0.36); \node at (2.6,1.5) {protein G};
  \node[below] at (2.6,1.1) {GDP $\to$ GTP};
  \draw[->, thick] (1.8,2.0) -- (2.2,1.75);
  % cyclase
  \draw[thick, fill=omDef!40] (4.2,1.7) rectangle (5.2,2.8); \node[above] at (4.7,2.85) {adenilil siklase};
  \draw[->, thick] (3.15,1.5) -- (4.2,1.85);
  % cAMP
  \node[omThm, anchor=west] at (5.35,1.55) {ATP $\to$ cAMP};
  \draw[->, thick, omThm] (4.7,1.7) -- (5.2,1.15);
  % PKA
  \draw[thick, fill=omThm!20] (7.3,0.9) ellipse (0.7 and 0.35); \node at (7.3,0.9) {kinase A};
  \draw[->, thick, omThm] (6.4,0.9) -- (6.6,0.9);
  % targets
  \draw[->, thick] (8.0,0.9) -- (9.0,0.9);
  \node[anchor=west, align=left] at (9.0,0.9) {fosforilase kinase $\to$ fosforilase\\$\to$ glikogen $\to$ glukosa};
  % amplification annotations
  \node[gray, align=center] at (1.4,0.3) {1 hormon};
  \node[gray, align=center] at (4.7,0.3) {$\sim 10^{2}$ protein G};
  \node[gray, align=center] at (7.3,0.3) {$\sim 10^{4}$ cAMP};
  \node[gray, align=center] at (10.2,0.3) {$\sim 10^{6}$ glukosa};
  % off switches
  \node[omExo!70!black, align=center, anchor=south] at (9.6,2.7) {padam: hidrolisis GTP (detik),\\fosfodiesterase memusnahkan cAMP,\\dan fosfatase membatalkan fosforilasinya};
\end{tikzpicture}

{\small Lintasan \omterm{prop:b2:cell-signalling:gpcr}{AMP siklik} dari adrenalin sampai glukosa. Tiap
tahapnya mengaktifkan banyak salinan tahap berikutnya: sebutir molekul
\omterm{def:b2:cell-signalling:kinds}{hormon} berakhir pada sejuta molekul hasil.}
\end{omfigure}

\begin{theorem}[Penguatan di dalam sebuah riam]\label{thm:b2:cell-signalling:amplification}
\index{penguatan (pensinyalan)}
Jika masing-masing dari $n$ tahap sebuah riam mengaktifkan $a$ molekul
tahap berikutnya selama masa hidup pengaktifannya sendiri, sebutir
molekul kurir menghasilkan $a^{n}$ molekul hasil akhirnya. Dengan empat
tahap penguat yang masing-masing seratus --- sebuah reseptor yang
mengaktifkan seratus \omterm{prop:b2:cell-signalling:gpcr}{protein G}, sebuah siklase yang membuat seratus
cAMP, sebuah kinase yang memfosforilasi seratus enzim, sebuah enzim
yang melepaskan seratus molekul glukosa --- $a^{n} = 10^{8}$; dalam
praktiknya sekitar sejuta, dalam beberapa detik, dari sebuah \omterm{def:b2:cell-signalling:kinds}{hormon}
pada $10^{-9}$ mol/L. Penguatan itulah sebabnya \omterm{def:b2:cell-signalling:kinds}{hormon} dapat mangkus
pada kepekatan sejuta kali lipat di bawah kepekatan metabolit yang ia
kendalikan, dan sebabnya riamnya harus dipadamkan pada tiap tahapnya:
sebuah penguat yang tidak dapat dimatikan adalah sebuah penyakit
(kolera, dan banyak kanker yang di dalamnya sebuah protein pensinyalan
terkunci menyala).
\end{theorem}
\begin{proof}
Tahap 1 menghasilkan $a$ molekul yang aktif; masing-masingnya
menghasilkan $a$ pada tahap 2, sehingga menjadi $a^{2}$; dengan
induksi, $a^{n}$ setelah $n$ tahap. Bati tiap tahapnya adalah jumlah
pergantian substrat yang dicapai sebuah katalis yang teraktifkan
sebelum ia dinonaktifkan: laju katalitiknya kali masa hidup aktifnya.
\end{proof}

\begin{proposition}[Reseptor tirosin kinase dan kalsium]\label{prop:b2:cell-signalling:rtk}
\index{reseptor tirosin kinase}\index{reseptor insulin}%
\index{pensinyalan kalsium}\index{kalmodulin}
Keluarga kedua reseptor permukaan --- bagi insulin, bagi faktor
pertumbuhan yang menggerakkan pembelahan sel, bagi isyarat pengimbas
perkembangannya --- adalah \emph{\omterm{prop:b2:cell-signalling:rtk}{reseptor tirosin kinase}}: protein
berlintasan tunggal yang, ketika ligannya menyatukan dua di antaranya,
saling memfosforilasi pada tirosinnya lalu menciptakan tapak
penambatan bagi sekumpulan protein adaptor. Dari sini berjalan
beberapa riam sekaligus: sebuah rantai tiga kinase (riam MAP kinase)
yang berakhir di dalam nukleusnya lalu menyalakan gen pertumbuhannya;
sebuah lipid kinase yang hasilnya merekrut kinase yang memerantarai
efek metabolis insulin --- memindahkan pengangkut glukosa ke
membrannya, menyalakan sintesis glikogennya --- dan lainnya. \omterm{prop:b2:cell-signalling:gpcr}{Kurir
kedua} semesta yang ketiga adalah \emph{kalsium}: ditahan di dalam
sitosolnya pada \qty{0.1}{\micro mol/L}, yaitu sepuluh ribu kali di
bawah kepekatan di luarnya, ia membanjir masuk lewat kanal yang dibuka
sebuah isyarat atau sebuah tegangan, atau keluar dari retikulum
endoplasmanya sebagai jawaban atas inositol trifosfat, lalu naik
menjadi mikromolar dalam hitungan milidetik; terikat pada
\emph{kalmodulin} ia mengaktifkan kinase dan pompa, dan terikat pada
troponin ia memicu kerutan (\cref{ch:b2:muscle-movement}). Pompa
mengembalikannya ke simpanannya dalam sedetik. Apa pun reseptornya,
bagian dalam selnya bertutur dalam beberapa bahasa --- cAMP, kalsium,
fosforilasi --- dan kekhususannya terletak pada protein mana yang
dimiliki tiap jenis sel untuk dijadikan sasarannya.
\end{proposition}

\begin{omfigure}
\includegraphics[width=0.6\linewidth]{images/book4/ai/b4-19-adrenaline-runner.jpg}

{\small Riamnya ketika bekerja. Sebuah ketakutan melepaskan adrenalin
dari medula adrenalnya; dalam hitungan detik \omterm{def:b2:heart:anatomy}{jantungnya} memacu,
hatinya mencurahkan glukosa, dan saluran napasnya melebar --- satu
kurir, beberapa reseptor, penguatan sejuta kali lipat di dalam tiap
selnya.}
\end{omfigure}

\section{Reseptor di dalamnya}

\begin{proposition}[Reseptor inti]\label{prop:b2:cell-signalling:nuclear}
\index{reseptor inti}\index{hormon steroid}
Steroid, tiroksin, asam retinoat, dan vitamin D berdifusi lewat
membrannya lalu mengikat reseptor di dalam sitosol atau nukleusnya yang
dirinya sendiri merupakan \emph{\omterm{prop:b2:cell-differentiation:transcriptional}{faktor transkripsi}}: kompleks
hormon--reseptornya mengikat runtunan DNA tertentu lalu menyalakan atau
memadamkan gen sasarannya. Tanggapannya lambat --- beberapa menit
sampai beberapa jam untuk membuat mRNA dan proteinnya --- dan bertahan
lama: kortisol mengimbas enzim glukoneogenesis selama berjam-jam,
estradiol membangun lapisan rahimnya selama berhari-hari, testosteron
membuat otot selama berminggu-minggu, tiroksin menetapkan laju
metabolisme tiap selnya. Beberapa steroid juga memiliki kerja yang
cepat lewat reseptor permukaan, dan beberapa lintasan permukaan juga
mencapai nukleusnya (riam MAP kinase), sehingga pembedaannya tidak
mutlak; tetapi kaidahnya tetap berlaku bahwa sebuah kurir yang larut
lipid mengubah \emph{apa} sebuah sel itu dengan mengubah gennya,
sedangkan kurir yang larut air mengubah apa yang sebuah sel
\emph{kerjakan} dengan mengubah enzimnya.
\end{proposition}

\begin{omfigure}
\begin{tikzpicture}[every node/.style={font=\scriptsize}, >=stealth]
  \foreach \x/\t in {0/{larut air: adrenalin, insulin, glukagon}, 7.4/{larut lipid: kortisol, estradiol, tiroksin}} \node[anchor=south, font=\scriptsize\bfseries] at (\x+2.6,3.4) {\t};
  % left: surface receptor
  \fill[omExo!30] (0,2.4) rectangle (5.2,2.8);
  \draw[thick, fill=omDef!40] (1.6,2.1) rectangle (2.3,3.1);
  \fill[omProp] (1.95,3.35) circle (0.1);
  \draw[->, thick] (1.95,2.1) -- (1.95,1.6) node[below, align=center] {kurir kedua,\\kinase: enzim\\disakelar dalam detik};
  \node[align=center] at (4.5,1.0) {mengubah apa yang\\dikerjakan selnya};
  % right: nuclear receptor
  \begin{scope}[xshift=7.4cm]
    \fill[omExo!30] (0,2.4) rectangle (5.2,2.8);
    \fill[omProp] (1.5,3.35) circle (0.1);
    \draw[->, omProp] (1.5,3.25) -- (1.5,1.9);
    \draw[thick, fill=omThm!20] (1.5,1.5) ellipse (0.72 and 0.3); \node at (1.5,1.5) {reseptor};
    \draw[thick] (0.6,0.4) -- (4.6,0.4); \draw[thick] (0.6,0.25) -- (4.6,0.25); \node[below] at (2.6,0.2) {DNA: gen sasaran nyala atau padam, berjam-jam};
    \draw[->, thick] (1.5,1.2) -- (2.4,0.5);
    \node[align=center] at (4.0,1.4) {mengubah apa\\selnya itu};
  \end{scope}
\end{tikzpicture}

{\small Dua rancangan. Sebuah kurir yang larut air bekerja di
permukaannya lalu mengubah keaktifan enzimnya dalam hitungan detik;
kurir yang larut lipid masuk, mengikat sebuah reseptor yang merupakan
\omterm{prop:b2:cell-differentiation:transcriptional}{faktor transkripsi}, lalu mengubah pengungkapan gennya selama
berjam-jam.}
\end{omfigure}

\section{Sebuah kasus terkerjakan: insulin, glukagon, dan titik acuan glukosanya}

\begin{proposition}[Dua hormon dan sebuah titik acuan]\label{prop:b2:cell-signalling:glucose}
\index{insulin}\index{glukagon}\index{glukosa darah}
\omterm{prop:b2:cell-signalling:glucose}{Glukosa darahnya} ditahan mendekati \qty{5}{mmol/L} (\qty{0.9}{g/L})
oleh dua \omterm{def:b2:cell-signalling:kinds}{hormon} berlawanan dari pulau pankreasnya. Setelah sebuah
santapan glukosanya yang naik diindra sel $\beta$ --- yang
memetabolismekannya, menutup sebuah kanal kalium ketika ATP-nya naik,
berdepolarisasi, memasukkan kalsium, lalu melepaskan \emph{insulin} ---
dan insulin, lewat \omterm{prop:b2:cell-signalling:rtk}{reseptor tirosin kinasenya}, membuat otot dan
lemaknya menyerap glukosa (dengan memindahkan pengangkut ke
membrannya), membuat hatinya menyimpannya sebagai glikogen dan lemak,
serta membuat seluruh jaringannya berhenti membakar lemak. Di antara
santapannya glukosanya yang turun membiarkan sel $\alpha$ melepaskan
\emph{glukagon}, yang, lewat cAMP di dalam hatinya, memecah glikogennya
lalu membuat glukosa baru dari asam amino; adrenalin mengerjakan hal
yang sama lebih cepat dalam keadaan darurat, dan kortisol serta \omterm{def:b2:cell-signalling:kinds}{hormon}
pertumbuhan selama berjam-jam. Tiap \omterm{def:b2:cell-signalling:kinds}{hormonnya} dilepaskan sebanding
dengan simpangan yang ia betulkan: sebuah umpan balik negatif dengan
dua efektor, satu bagi tiap arahnya, dan sebuah titik acuan yang dibela
sampai seperlima. Otaknya, yang membakar \qty{120}{g} glukosa tiap hari
dan tidak menyimpan apa pun, adalah organ yang dilindungi gelungnya;
diabetes tipe 1 adalah hilangnya sel $\beta$-nya, dan perjalanannya
yang tak terobati --- glukosa pada \qty{20}{mmol/L}, lemak yang dibakar
karena kekurangan insulin, asam di dalam \omterm{def:b2:blood-circulation:blood}{darahnya} --- adalah rupa
gelungnya dengan satu lengan yang disingkirkan.
\end{proposition}

\begin{omfigure}
\begin{tikzpicture}
\begin{axis}[
  width=0.82\linewidth, height=5.6cm,
  axis x line=bottom, axis y line=left,
  xmin=0, xmax=6, ymin=0, ymax=1.1,
  xlabel={jam setelah sebuah santapan}, ylabel={nisbi terhadap puncaknya},
  xlabel style={font=\small, at={(axis description cs:0.5,-0.12)}, anchor=north},
  ylabel style={font=\small},
  tick label style={font=\small}, xtick={0,1,2,3,4,5,6}, ytick={0,0.5,1},
  legend style={font=\scriptsize, at={(0.5,-0.28)}, anchor=north, draw=none, legend columns=3},
  clip=false,
]
\addplot[very thick, omDef, domain=0:6, samples=80] {0.55 + 0.45*exp(-((x-0.9)^2)/0.5) - 0.08*exp(-((x-4)^2)/2)};
\addlegendentry{glukosa darah (5 mmol/L pada 0.55)}
\addplot[very thick, omThm, domain=0:6, samples=80] {0.15 + 0.85*exp(-((x-1.1)^2)/0.45)};
\addlegendentry{insulin}
\addplot[very thick, omProp, domain=0:6, samples=80] {0.75 - 0.5*exp(-((x-1.2)^2)/0.8) + 0.25*(1-exp(-x/5))*0.6};
\addlegendentry{glukagon}
\end{axis}
\end{tikzpicture}

{\small Glukosa, insulin, dan glukagon setelah sebuah santapan.
Insulinnya naik bersama glukosanya lalu menurunkannya; ketika
glukosanya turun menuju titik acuannya, glukagonnya naik untuk
menahannya di sana.}
\end{omfigure}

\begin{example}[Kecepatan, jangkauan, dan ongkos kedua sistemnya]\label{ex:b2:cell-signalling:compare}
Adrenalin di dalam \omterm{def:b2:blood-circulation:blood}{darahnya} mencapai tiap sel dalam kira-kira semenit
(yaitu waktu peredarannya) dan bekerja selama beberapa menit; sebuah
impuls saraf mencapai sasarannya dalam hitungan milidetik dan bekerja
selama beberapa milidetik. Sebuah \omterm{def:b2:cell-signalling:kinds}{hormon} nyaris tidak berongkos apa pun
tiap sel yang dicapainya --- satu nanomol yang tersebar ke seluruh
tubuhnya --- tetapi tidak dapat mengalamati satu sel; sebuah saraf
mengalamati satu sel secara tepat tetapi harus membangun dan merawat
sebuah kabel menujunya. Tubuhnya memakai keduanya: saraf simpatis pada
\cref{ch:b2:blood-pressure} bekerja pada \omterm{def:b2:heart:anatomy}{jantung} dan \omterm{def:b2:blood-circulation:vessels}{arteriolnya} dalam
sedetik, dan medula adrenalnya, yang dirinya sendiri merupakan sebuah
ganglion simpatis yang terubah, menyokongnya dengan mencurahkan molekul
yang sama ke dalam \omterm{def:b2:blood-circulation:blood}{darahnya} bagi seluruh tubuhnya. Dan pada sinapsisnya
kedua rancangannya berjumpa, karena sebuah \omterm{def:b2:cell-signalling:kinds}{neurotransmiter} adalah
sebuah kurir yang bekerja pada sebuah reseptor lewat mekanisme yang
sama --- sebuah kanal yang membuka, sebuah \omterm{prop:b2:cell-signalling:gpcr}{protein G} yang menyakelar
--- hanya saja melintasi dua puluh nanometer alih-alih satu meter.
\end{example}

\section{Latihan}

\begin{exercise}[$\star$]\label{exo:b2:cell-signalling:1}
Golongkanlah sebagai endokrin, parakrin, autokrin, atau sinaptik, dan
sebagai larut air atau larut lipid: insulin, kortisol, asetilkolin pada
sebuah sinapsis, histamin pada sebuah luka, estradiol, adrenalin dari
adrenalnya, nitrogen monoksida dari endoteliumnya.
\end{exercise}

\begin{exercise}[$\star$]\label{exo:b2:cell-signalling:2}
Sebutkan langkahnya dari adrenalin pada membran sel hatinya sampai
glukosa di dalam \omterm{def:b2:blood-circulation:blood}{darahnya}, dan sebutkan sakelar pemadam tiap
langkahnya.
\end{exercise}

\begin{exercise}[$\star$]\label{exo:b2:cell-signalling:3}
Bandingkan reseptor permukaan dan \omterm{prop:b2:cell-signalling:nuclear}{reseptor inti}: apa yang mengikatnya,
di mana letaknya, seberapa cepat ia bekerja, dan apa yang ia ubah.
\end{exercise}

\begin{exercise}[$\star$]\label{exo:b2:cell-signalling:4}
Apakah ketiga \omterm{prop:b2:cell-signalling:gpcr}{kurir kedua} semesta yang disebutkan pada bab ini,
bagaimana masing-masingnya dinaikkan, dan bagaimana masing-masingnya
disingkirkan?
\end{exercise}

\begin{exercise}[$\star\star$]\label{exo:b2:cell-signalling:5}
Sebuah reseptor memiliki $K_d = \qty{2}{nmol/L}$. Hitunglah
keterisiannya pada 0.2, 2, 20, dan \qty{200}{nmol/L}. Sebuah sel
melipatduakan jumlah reseptornya: apa yang berubah, keterisiannya atau
jumlah reseptor yang terisi?
\end{exercise}

\begin{exercise}[$\star\star$]\label{exo:b2:cell-signalling:6}
Sebuah riam memiliki tahap dengan bati 50, 200, 100, dan 300.
Hitunglah penguatan keseluruhannya. Jika \omterm{def:b2:cell-signalling:kinds}{hormonnya} berada pada
\qty{1}{nmol/L} dan \qty{1}{\%} reseptornya ($10^{4}$ tiap sel) terisi,
berapa banyak molekul hasil yang dibuat selnya tiap detik jika bati
tiap tahapnya adalah bati tiap detik?
\end{exercise}

\begin{exercise}[$\star\star$]\label{exo:b2:cell-signalling:7}
Toksin kolera mengunci \omterm{prop:b2:cell-signalling:gpcr}{protein G} perangsangnya dalam keadaan terikat
GTP di dalam sel ususnya. Jelaskan, langkah demi langkah, mengapa
pasiennya kehilangan berliter-liter cairan tiap hari, dan mengapa
efeknya bertahan berhari-hari setelah bakterinya lenyap.
\end{exercise}

\begin{exercise}[$\star\star$]\label{exo:b2:cell-signalling:8}
Kalsium sitosolnya naik dari \qty{0.1}{} menjadi \qty{1}{\micro mol/L}
di dalam sebuah sel \qty{2000}{\micro m^3} ketika sebuah kanal membuka.
Berapa banyak ion kalsium yang masuk? Jika sebuah kanal terbuka tunggal
melewatkan $10^{6}$ ion tiap detik, berapa lama satu kanal
memerlukannya, dan mengapa sebuah sel membuka banyak kanal sekaligus?
\end{exercise}

\begin{exercise}[$\star\star$]\label{exo:b2:cell-signalling:9}
Insulin dan glukagon keduanya bekerja pada hatinya, yang satu lewat
sebuah tirosin kinase dan yang lain lewat cAMP. Jelaskan bagaimana sel
hatinya memadukan kedua isyaratnya, apa yang terjadi pada glikogennya
ketika keduanya tinggi, dan mengapa nisbahnya lebih penting daripada
salah satu arasnya.
\end{exercise}

\begin{exercise}[$\star\star\star$]\label{exo:b2:cell-signalling:10}
Adrenalin menaikkan cAMP di dalam sebuah sel hati dalam \qty{5}{s} dan
keluaran glukosanya dalam \qty{30}{s}; kortisol menaikkan enzim
glukoneogenik hatinya selama \qty{4}{h}. Dengan memakai difusi,
transkripsi (\qty{50}{} nukleotida tiap detik pada sebuah gen
\qty{3000}{} nukleotida), translasi (\qty{5}{} asam amino tiap detik
pada sebuah protein \num{500} residu), dan keperluan menimbun beribu
molekul enzim, terangkanlah kedua skala waktunya.
\end{exercise}

\begin{exercise}[$\star\star\star$]\label{exo:b2:cell-signalling:11}
Sebuah \omterm{def:b2:genome-diversification:mutation}{mutasi} membuat sebuah \omterm{prop:b2:cell-signalling:rtk}{reseptor tirosin kinase} faktor
pertumbuhan berdimer, dan dengan demikian memberi isyarat, tanpa
ligannya. Ramalkan perilaku selnya, jelaskan mengapa ini adalah langkah
yang lazim pada kanker, dan katakan apa yang akan dikerjakan sebuah
obat yang menghalangi tapak ATP kinasenya.
\end{exercise}

\begin{exercise}[$\star\star\star$]\label{exo:b2:cell-signalling:12}
``Sistem endokrin menyiarkan; sistem saraf mengabel.'' Bahaslah kedua
rancangannya dalam hal kecepatan, jangkauan, kekhususan, dan ongkosnya,
lalu katakan di mana masing-masingnya tak tergantikan dan di mana
tubuhnya memakai keduanya bagi pesan yang sama.
\end{exercise}

\section{Soal: Sebuah Hormon dari Darah sampai Gen}

\begin{problem}[{Soal akhir pekan --- adrenalin diikuti dari kelenjar
adrenal sampai glukosa sebuah sel hati, dengan kepekatan, keterisian
reseptor, bati riam, dan waktunya dihitung, lalu kortisol diikuti
sampai sebuah gen, dan gelung glukosanya dianalisis, berakhir pada
keterisiannya, penguatannya, dan waktu yang diperlukan tiap hormonnya}]
\label{pb:b2:cell-signalling:1}
Datanya: medula adrenalnya melepaskan \qty{1}{nmol} adrenalin ke dalam
\qty{5}{L} \omterm{def:b2:blood-circulation:blood}{darah} ketika ketakutan; $K_d$ adrenalin pada reseptor
hatinya adalah \qty{1}{nmol/L}; sebuah sel hati (\qty{5000}{\micro
m^3}) membawa $2\times 10^{4}$ reseptor. Bati riamnya tiap detik
pengaktifan: reseptor $\to$ \omterm{prop:b2:cell-signalling:gpcr}{protein G} 20, \omterm{prop:b2:cell-signalling:gpcr}{protein G} $\to$ cAMP 100
(yaitu pergantian siklasenya selama masa hidup \omterm{prop:b2:cell-signalling:gpcr}{protein G} yang 3 detik),
cAMP $\to$ kinase A 1 (stoikiometrik), kinase A $\to$ fosforilase
kinase 50, $\to$ fosforilase 50, $\to$ glukosa 1000. Hatinya menahan
\qty{100}{g} glikogen. Kortisol: transkripsi pada \qty{50}{} nukleotida
tiap detik pada sebuah gen \qty{4000}{} nukleotida, translasi pada
\qty{5}{} residu tiap detik pada sebuah enzim \num{600} residu, \num{20}
ribosom tiap mRNA, \num{100} mRNA dibuat tiap jam, dan sebuah efek yang
memerlukan $10^{5}$ molekul enzim.

\medskip
\noindent\textbf{Bagian I --- \omterm{def:b2:cell-signalling:kinds}{Hormonnya} tiba.}
\begin{enumerate}
  \item Hitunglah kepekatan adrenalin di dalam \omterm{def:b2:blood-circulation:blood}{plasmanya} setelah
        pelepasannya, dengan mengabaikan penyingkirannya.
  \item Hitunglah \omterm{thm:b2:cell-signalling:occupancy}{keterisian reseptornya} dan jumlah reseptor yang
        terisi tiap sel hatinya.
  \item Adrenalin disingkirkan dengan paruh waktu \qty{2}{min}.
        Berapakah kepekatan dan keterisiannya setelah \qty{6}{min}?
  \item Berapa lama adrenalinnya memerlukan waktu untuk mencapai
        hatinya dari adrenalnya, mengingat waktu peredarannya
        kira-kira semenit? Mengapa saraf simpatis menuju hatinya
        bekerja lebih cepat?
  \item Sebuah sel dengan $2\times 10^{5}$ reseptor dan $K_d$ yang
        sama: apa yang berubah?
  \item Mengapa $K_d$ diserasikan dengan kepekatan yang beredar ---
        apa yang akan dikerjakan sebuah reseptor dengan
        $K_d = \qty{1}{\micro mol/L}$ terhadap sebuah \omterm{def:b2:cell-signalling:kinds}{hormon} nanomolar?
\end{enumerate}

\medskip
\noindent\textbf{Bagian II --- Riamnya.}
\begin{enumerate}[resume]
  \item Hitunglah bati keseluruhannya dari satu reseptor yang terisi
        sampai molekul glukosa tiap detik.
  \item Hitunglah molekul cAMP yang dibuat selnya pada detik pertama,
        dan kepekatan yang dicapainya di dalam volumenya.
  \item Hitunglah keluaran glukosa satu sel hati tiap detik pada
        keterisian pada pertanyaan 2, dan tiap menit.
  \item Hatinya memiliki $2\times 10^{11}$ sel. Hitunglah keluaran
        yang diramalkan batinya, dalam gram tiap menit, lalu
        bandingkan dengan \qty{5}{g/min} yang sungguh teramati. Apa
        yang membatasi keluaran yang sebenarnya?
  \item Berapa lama \qty{100}{g} glikogennya akan bertahan pada laju
        ini? Mengapa ia tidak habis ketika ketakutan?
  \item Sebutkan sakelar pemadam pada tiap tahap penguatnya dan
        katakan mana yang bekerja paling cepat.
  \item Sebuah penghambat fosfodiesterase (kafeina) hadir. Ramalkan
        efeknya pada riamnya dan pada keluaran glukosanya.
\end{enumerate}

\medskip
\noindent\textbf{Bagian III --- Kortisol menuju sebuah gen.}
\begin{enumerate}[resume]
  \item Hitunglah waktu untuk mentranskripsi satu mRNA.
  \item Hitunglah waktu untuk mentranslasi satu molekul enzim, dan
        jumlah molekul enzim tiap mRNA tiap jam (tiap mRNA hidup
        sejam, dengan ribosom yang mulai satu demi satu).
  \item Hitunglah molekul enzim yang dibuat tiap jam dari \num{100}
        mRNA-nya, dan waktu untuk mencapai $10^{5}$.
  \item Tambahkan waktu bagi kortisol untuk masuk, mengikat, dan
        mencapai DNA-nya (beberapa menit) lalu terangkanlah
        \qty{4}{h} yang teramati.
  \item Mengapa tubuhnya tidak akan memakai sebuah riam seperti
        riam adrenalin bagi tugas kortisolnya, dan mengapa tidak
        memakai mekanisme kortisol bagi sebuah ketakutan?
  \item Kortisol juga menaikkan jumlah reseptor adrenalin di dalam sel
        hatinya. Jelaskan bagaimana sebuah \omterm{def:b2:cell-signalling:kinds}{hormon} yang lambat dapat
        menguatkan \omterm{def:b2:cell-signalling:kinds}{hormon} yang cepat, dan apa yang dikerjakan hal ini
        terhadap kurva keterisian \omterm{def:b2:cell-signalling:kinds}{hormon} yang cepat itu.
\end{enumerate}

\medskip
\noindent\textbf{Bagian IV --- Gelungnya.}
\begin{enumerate}[resume]
  \item Setelah sebuah santapan glukosanya naik dari 5 menjadi
        \qty{8}{mmol/L}. Perikanlah tanggapan sel $\beta$-nya dari
        glukosa sampai pelepasan insulin, dengan menyebutkan
        gandengannya.
  \item Insulin bekerja lewat sebuah tirosin kinase, glukagon lewat
        cAMP. Jelaskan mengapa satu sel dapat menuruti keduanya dan
        bagaimana enzim glikogen hatinya ditetapkan oleh nisbahnya.
  \item Seorang pasien tidak memiliki sel $\beta$. Ramalkan glukosanya
        ketika berpuasa dan setelah santapan, lalu jelaskan mengapa
        lemaknya dibakar dan asam muncul.
  \item Insulin yang disuntikkan tidak memiliki umpan balik. Jelaskan
        bahaya sebuah lari yang luar biasa panjang dengan takaran yang
        biasa.
  \item Bandingkan gelung glukosanya dengan \omterm{def:b2:blood-pressure:baroreflex}{barorefleks} pada
        \cref{ch:b2:blood-pressure}: pengindra, efektor, bati,
        kecepatan.
  \item Nyatakan hasilnya: keterisian adrenalin setelah pelepasannya,
        bati riamnya, keluaran glukosa hatinya, dan waktu yang
        diperlukan kortisol.
\end{enumerate}
\end{problem}
